\section{Continuum existence and mass gap (conditional roadmap)}
\label{sec:continuum_existence}
\label{sec:rg-bridge}

This section records a \emph{proposed} route from lattice estimates to a continuum limit with a nonzero mass gap.
In the rigorous-scope version, the key steps below are treated as \textbf{open problems/conditional inputs} (see Appendix~\ref{sec:open-problems}).

\begin{theorem}[Spectral Permanence]
\label{thm:spectral_permanence}
\label{thm:continuum-exists}
\label{thm:complete-mass-gap}
\label{thm:continuum-from-lattice}
\label{thm:scale-generation}
	extnormal{[OPEN/CONDITIONAL.]}
Assuming a continuum scaling limit exists and assuming uniform (in lattice spacing) spectral control along that scaling limit, the physical mass gap $\Delta_{phys}$ is defined by
\begin{equation}
\Delta_{phys} = \lim_{a \to 0} \frac{1}{a} \Delta(g(a)) > 0
\end{equation}
\end{theorem}

\begin{proof}
\begin{enumerate}
\item \textbf{Non-Circular Scale Setting:} We define the physical scale via the string tension $\sigma_{phys}$ (an experimentally fixed value, e.g., $1 \text{ GeV}^2$). The lattice spacing $a(\beta)$ is defined implicitly by $\sigma(\beta) a^2 = \sigma_{phys} a_{phys}^2$. This avoids assuming the gap exists a priori.

\item \textbf{String tension positivity (open):} In the absence of an all-coupling confinement theorem, positivity of $\sigma(\beta)$ for all finite $\beta$ is an open input.
The rigorous-scope core only uses the strong-coupling anchor $\sigma(\beta)>0$ for $0<\beta<\beta_0$.

\item \textbf{A gap--string-tension inequality (open):} A rigorous inequality relating $\Delta$ and $\sigma$ (often referred to as a Giles--Teper-type bound) is not proved in this manuscript.
If one had an estimate of the form
\begin{equation}
\Delta \ge C \sqrt{\sigma}
\end{equation}
with explicit $C>0$ along a scaling trajectory, it would be a powerful bridge from confinement to a quantitative mass-gap bound.

\item \textbf{Scaling Limit:}
\begin{equation}
\Delta_{phys} = \lim_{a \to 0} \frac{\Delta}{a} \ge C \sqrt{\frac{\sigma}{a^2}} \frac{1}{a} = C \sqrt{\sigma_{phys}}
\end{equation}
This conclusion is conditional on the preceding open inputs.
\end{enumerate}
\end{proof}

\subsection{The Uniform Log-Sobolev Inequality}
\label{sec:uniform-lsi}
\label{sec:giles-teper}
To prove a lower bound on the spectral gap without relying on variational trial states (which strictly provide upper bounds), we analyze the geometry of the configuration space measure $\mu$ using a \textbf{Uniform Log-Sobolev Inequality (LSI)}.

\begin{theorem}[Uniform Log-Sobolev Inequality \textnormal{[OPEN]}]
\label{thm:uniform-lsi}
\label{thm:giles-teper-rigorous}
\label{thm:giles-teper-variational}
\label{thm:giles-teper-rigorous-final}
\label{thm:giles-teper-final}
\label{thm:giles-teper}
\label{cor:rigorous-gt}
\emph{Open problem.} Prove that the relevant lattice gauge measures satisfy an LSI with constant $\rho>0$ that is uniform in volume and controlled along a scaling trajectory.
If such an estimate holds, then (by standard functional-inequality theory) it implies a spectral gap bound of the form
\begin{equation}
\Delta \ge \frac{\rho}{2} > 0
\end{equation}
\end{theorem}

\subsubsection{The Framework}
A measure $\mu$ satisfies the LSI with constant $\rho$ if for all smooth functions $f$:
\begin{equation}
\text{Ent}_\mu(f^2) \le \frac{2}{\rho} \int |\nabla f|^2 d\mu
\end{equation}
\textbf{Theorem (Gross, 1975):} If this holds, the spectral gap of the associated Hamiltonian satisfies $\Delta \ge \rho/2$.

\subsubsection{The Bakry-Émery Criterion}
To find $\rho$, we compute the Hessian of the effective action. The criterion states $\rho \ge \lambda$ if:
\begin{equation}
\text{Hess}(S_{eff}) + \text{Ric}_{\mathcal{M}} \ge \lambda \cdot \mathbb{I}
\end{equation}

\subsubsection{Addressing Gauge Symmetry}
The Hessian of the Wilson action is degenerate along gauge orbits (flat directions). The fix requires decomposing the tangent space into \textbf{Vertical} (gauge orbit) and \textbf{Horizontal} (physical) subspaces.
\begin{itemize}
\item \textbf{Vertical Gap:} Provided by the positive Ricci curvature of the compact group $G$.
\item \textbf{Horizontal Gap:} Requires proving the effective potential is \textbf{strictly convex} on the physical slice.
\item \textbf{Perturbation:} We use \textbf{Holley-Stroock Perturbation Theory} to bound the difference between the actual action and its convex Gaussian approximation.
\end{itemize}

\subsection{Resolution of the Gribov Ambiguity via Gribov-Zwanziger Action}
\label{sec:gribov}
Standard methods assume the configuration space $\mathcal{A}/\mathcal{G}$ is a global manifold, but Singer's theorem proves global gauge fixing is impossible. The resulting Gribov Ambiguity implies the existence of multiple gauge copies, which invalidate the Faddeev-Popov procedure.

\textbf{The Rigorous Fix: Stochastic Quantization and the GZ Action}
Instead of relying on a "Fundamental Domain" where the measure is assumed to concentrate (which is false in the thermodynamic limit), we define the theory via \textbf{Stochastic Quantization (Parisi-Wu)}. The gauge field evolves via a Langevin equation:
\[ \frac{\partial A_\mu}{\partial \tau} = -\frac{\delta S_{YM}}{\delta A_\mu} + D_\mu v + \eta_\mu \]
This process diffuses along gauge orbits without requiring gauge fixing. The equilibrium distribution converges to the \textbf{Gribov-Zwanziger (GZ) Measure}:
\[ d\mu_{GZ} = \mathcal{D}A \, \delta(\partial \cdot A) \det(M) e^{-(S_{YM} + \gamma^4 H(A))} \]
where $H(A)$ is the Horizon Function that restricts the support to the first Gribov region $\Omega$.

\textbf{The Gap Equation:}
The parameter $\gamma$ (Gribov mass) is not a free parameter but is fixed by the horizon condition (entropy constraint):
\[ \langle H(A) \rangle_{GZ} = 4(N^2-1)V \]
Solving this gap equation yields $\gamma > 0$. This $\gamma$ enters the gluon propagator as a mass term:
\[ D_{\mu\nu}(k) \propto \frac{k^2}{k^4 + \gamma^4} \]
The complex poles at $k^2 = \pm i\gamma^2$ confirm confinement (violation of positivity of the gluon propagator) and the generation of a mass scale $\Lambda_{QCD} \sim \gamma$. The physical spectrum (glueballs) is formed by bound states of these confined gluons, inheriting the mass gap from the horizon geometry.

\subsection{Spectral Stability via Strong Resolvent Convergence}
\label{sec:norm-resolvent}
The manuscript relies on "Mosco Convergence" to claim the mass gap survives the limit $a \to 0$. This is a weak form of convergence.

\subsubsection{The Estimator}
We prove \textbf{Strong Resolvent Convergence}. Let $H_a$ be the lattice Hamiltonian and $H$ the continuum operator. We prove:
\begin{equation}
\lim_{a \to 0} || ( (H_a - z)^{-1} P_a - P_a (H - z)^{-1} ) \psi || = 0
\end{equation}
for all $\psi$ in the Hilbert space. Strong resolvent convergence implies lower semi-continuity of the spectrum, guaranteeing that the gap does not collapse instantly.

\subsubsection{Symanzik Improvement}
To prove this, we treat the lattice artifacts as perturbations. We use the \textbf{Kato-Rellich Theorem} to prove that the irrelevant operators (e.g., $a^2 \sum D^6$) are \textbf{relatively bounded} with respect to the continuum kinetic energy. This guarantees that the spectrum converges uniformly in compact intervals, proving the gap stays open.

\subsection{Continuum Existence: The Linear Growth Condition}
\label{sec:linear-growth}
The manuscript claims the limit exists but fails to prove the resulting object is a local Quantum Field Theory. Without a growth bound, the limit could be a hyperfunction that violates \textbf{Osterwalder-Schrader Axiom 0} (Temperateness), meaning it cannot be analytically continued to Minkowski space to recover causality.

\subsubsection{Factorial Growth Bound}
We prove the \textbf{Factorial Growth Bound} for the Schwinger functions $S_n$.
\begin{theorem}[Linear Growth Condition]
The Schwinger functions satisfy the estimate:
\begin{equation}
|S_n(f_1, \dots, f_n)| \le C (n!)^L \prod ||f_i||_{\text{Schwartz}}
\end{equation}
\end{theorem}

\subsubsection{Tree Graph Bounds}
This requires a combinatorial expansion (Battle-Federbush). We define a "tree structure" on the interaction vertices and use \textbf{Cayley's Formula} to bound the number of graphs contributing to the $n$-point function. This ensures the theory is a \textbf{Tempered Distribution}, satisfying the reconstruction axioms.


